\documentclass[11pt]{article}
\PassOptionsToPackage{numbers}{natbib}
\usepackage[review]{acl}
\makeatletter\acl@anonymizefalse\makeatother
\usepackage{times}
\usepackage{latexsym}
\usepackage[T1]{fontenc}
\usepackage[utf8]{inputenc}
\usepackage{microtype}
\usepackage[htt]{hyphenat}
\usepackage{booktabs}
\usepackage{amsmath}
\usepackage{graphicx}
\sloppy
\setlength{\tabcolsep}{4pt}
\title{Team Name at SemEval-2027 Task 1: Multi-evidence re-ranking with a frozen first stage}
\author{Author Name \\ Affiliation \\ \texttt{email@example.org}}
\date{}
\begin{document}
\maketitle
\begin{abstract}
RETECO Track 1a asks for documents that answer a \emph{complete} question about a historical period: long queries whose temporal and exclusion conditions sit in the body rather than the title, and judgements that mark \textbf{several} documents per question (3.44 on average, over \textbf{2.34 annotated steps}). We describe a two-stage ranker - a zero-shot dense retriever plus a windowed cross-encoder fine-tuned on the 1,071 training queries that carry at least one positive among their dense top-100 candidates - that scores 0.3581 domain-macro nDCG@10 on the 13-domain dev set, against 0.2943 for the dense stage and 0.1147 for the organisers’ BM25 baseline. Against the same pipeline with the cross-encoder left at zero-shot (0.3249), the fine-tuned reranker contributes +0.0332 (domain-clustered 95\% CI [+0.0192, +0.0472]) at a fusion form fixed before fine-tuning, and reproduces at +0.0326 / +0.0279 under two further seeds. The remaining headroom is only partly recall: at the 100 candidates the reranker sees, 80.7\% of annotated steps have at least one relevant document (96.3\% within the dense top-1000), yet the fused top ten covers only 59.5\% of the annotated steps and the cross-encoder’s own top ten 51.4\%. Ten temporal probes, one strictly pre-registered test (+0.0039 against a +0.02 bar) and a fixed year-consistency reordering that \emph{lowers} nDCG@10 in 7 of 8 domains do not close it. Supervising every annotated step to have a positive ranked high is a wash (0.3593 and 0.3600 against 0.3581, both intervals include zero) and leaves coverage where it was. Restoring the 140 training queries with an empty candidate pool is a wash too (0.3565). Only the constraint-faithful 64-word query pays a price: 0.3443, 0.0138 below the frozen system. Every intervention in the six-step ladder is neutral or negative on dev, and the frozen system is kept.

\end{abstract}
\section*{1 Introduction}
Track 1a gives 13 independent corpora, long queries, binary qrels and no timestamps. Our entry has two parts: a strong, reproducible zero-shot-to-supervised pipeline, and a diagnosis of why the obvious temporal fixes do not work.

Contributions: (i) a baseline ladder over all 13 domains under the official domain-macro statistic, with paired bootstraps over queries and over domains; (ii) a supervised cross-encoder that adds +0.0332 without tuning the fusion weight on dev; (iii) a multi-evidence analysis showing that the system’s failure mode is "one easy document" inside a candidate pool that is itself incomplete; (iv) negative evidence on explicit temporal modelling: one pre-registered test and a fixed year-consistency rule that fails; (v) a decomposition of why the cross-encoder with the best standalone score is not the one that fuses best - anchoring, attenuation and score calibration - and what that implies for using a single-stage metric to select checkpoints.

\emph{Relation to prior work.} The reranker is the standard retrieve-then-rerank cross-encoder \citep{nogueira2019passage}, initialised from M3-Embedding \citep{chen2024m3}. Our multi-evidence objective is a sibling of listwise and LambdaRank training \citep{burges2010lambdarank,wang2018lambdaloss}, with one difference: instead of ordering one list, it asks that every official step of a question own a document in the top ten. Multi-stage fusion has its own line - list-aware reranking of the first stage’s list \citep{zhang2022hlatr} and joint selection of documents that cover the answer set \citep{mao2021jpr} - but both learn the second stage jointly with the list it consumes (HLATR adds a list-aware module over retrieved and reranked lists; JPR reranks the set autoregressively), whereas we freeze the first stage and ask what one second stage can add at a fixed fusion weight. The contribution is the measured behaviour of that frozen system, not a new fusion architecture. Temporal retrieval is the benchmark’s own subject \citep{abdallah2026tempo}; our contribution there is negative evidence about what does not transfer from the two-domain pilot to the full suite.

\section*{2 Task, metric, data}
\textbf{Metric.} Official aggregation is domain-macro nDCG@10: per-query nDCG@10, averaged within domain, then over the 13 domains, so a 3-query domain weighs as much as a 240-query one, and queries missing from a run score zero. dev = 519 queries. 26.6\% of dev questions receive exactly zero from the frozen system, and 37.7\% of those have no relevant document in the top 100 at all. Those zeros dominate the statistic: had the 138 zero questions merely reached half their domain’s average, domain-macro would be 0.4283 instead of 0.3581 (full average: 0.4985, \texttt{results/ZERO\_HEADROOM.json}).

\textbf{Data.} 1,167,159 documents, 1,211 train queries, binary qrels (so nDCG@10 is discounted precision at 10). Per-domain counts in \texttt{results/corpus\_manifest.json}. Queries carry a title and a body separated by \texttt{<p>}; the body holds the temporal range, exclusions and comparisons (Section 5.3).

\section*{3 System}
\textbf{Dense first stage.} \texttt{mxbai-embed-large-v1}, full query with the model’s retrieval instruction, mean pooling, cosine similarity, fp16, \texttt{top-100} candidates per query. Dev: 0.2943 (0.2956 with CLS pooling: a wash).

\textbf{Windowed cross-encoder.} \texttt{bge-reranker-v2-m3} (568M). Documents are split into 250-token windows with stride 192, each window is scored against the query, and the document score is the maximum over its windows, i.e. a document is as relevant as its best passage; query capped at 256 tokens. Pooling was fixed a priori and is not ablated (Section 7). Zero-shot at fusion 1:1: 0.3249.

\textbf{Fusion.} Per-query z-normalisation of the dense and cross-encoder lists, combined at 1:1. The weight was fixed on the training split before any dev inspection; Section 5.1 shows the choice is not what produces the gain.

\textbf{Supervised fine-tuning.} Listwise softmax over a 20-document group sampled from each training query’s 100 candidates - 21,420 query-document pairs in 1,071 groups, the log-sum-exp running over the group’s positives - 8 of 24 layers frozen, lr 1.5e-5, seed 17, 2 epochs; epoch 1 was kept, by a 0.3029-against- 0.3012 internal-validation margin. Checkpoint selection used the cross-encoder’s own ordering on 45 held-out training queries, not the fused system; that run predates domain-macro logging and selected on the plain mean, and on every later arm the two criteria pick the same epoch. Frozen dev: \textbf{0.3581}.

\textbf{Multi-evidence (step) supervision.} Track 1a gives one relevance judgement per query, but the organisers’ qrels also carry \emph{step} annotations: a question that needs several pieces of evidence is decomposed into steps, and each step lists the documents that satisfy it. We use them only in training; at test time the ranker still sees the full question and one document. With \texttt{P\_j} the step-\texttt{j} positives inside a query’s candidate list, the ordinary listwise term \texttt{-log sum\_\{d in P\} p(d)} is dominated by the single easiest positive, so we add

\texttt{`} L\_step = -(1/m) * sum\_j log sum\_\{d in P\_j\} p(d),    L = L\_list + lambda * L\_step \texttt{`}

which forces \emph{every} step to have a positive ranked high instead of letting one easy document carry the query. Steps whose positives are absent from the candidate list are skipped (a step is covered for 91.2\% of step instances under our candidate construction); \texttt{m} counts the covered steps. \texttt{lambda} is not tuned on dev: two values (0.3, 1.0) are trained on the same groups, and both raise the train-internal split above the baseline (0.2882 and 0.2913 against 0.2865), yet both are rejected on dev (S5.6).

\section*{4 Results}
\begin{table}[t]
\centering
\footnotesize
\resizebox{\columnwidth}{!}{%
\begin{tabular}{lrrr}
\toprule
system & dev domain-macro nDCG@10 & $\Delta$ vs zero-shot z1:1 & domain-clustered 95\% CI \\
\midrule
official BM25 & 0.1147 & --- & --- \\
dense (mxbai, top-100) & 0.2943 & --- & --- \\
dense + zero-shot CE, z 1:1 & 0.3249 & --- & --- \\
+ SFT CE, 1 epoch & 0.3336 & +0.0087 & [+0.0037, +0.0142] \\
\textbf{+ SFT CE, 2 epochs (frozen)} & \textbf{0.3581} & \textbf{+0.0332} & \textbf{[+0.0192, +0.0472]} \\
same recipe, seed 23 & 0.3575 & +0.0326 & [+0.0184, +0.0470] \\
same recipe, seed 29 (two independent reruns) & 0.3528 / 0.3527 & +0.0279 & [+0.0160, +0.0405] \\
\bottomrule
\end{tabular}}%
\end{table}

The gain holds when the three smallest domains are dropped (0.3267 against 0.2985 on domains with n >= 19 queries) and in 11 of 13 domains; the two exceptions, bitcoin and quant, are flat (-0.0023 and -0.0009). Seed spread over three seeds is 0.0029, an order of magnitude below the effect, while a fourth run of the same recipe through the current code path scores 0.3575: the retraining band is 0.3527-0.3581 and the frozen artifact sits at its top (S5.5). Rank-fusion forms were screened too and are reported in S5.1, where the RRF list that out-of-fold selection prefers is worse than the z 1:1 fusion used here.

\section*{5 Analysis}
\subsection*{5.1 Is the gain an artefact of the fusion weight?}
No. A sweep over nine weights (0.5 to 2 on the cross-encoder, z-normalised) is positive everywhere (+0.0089 to +0.0346; the w=1 row reproduces the headline within bootstrap noise), and refitting the weight out-of-fold on 45 held-out training queries selects a dense-heavy RRF 1:4 that is \emph{worse} than z 1:1, so the gain is not bought by tuning on dev.

\subsection*{5.2 The bottleneck is multi-evidence coverage over an incomplete pool}
We measured, per system and without extra annotation, how much of the official step structure reaches the top ten (\texttt{scripts/dev\_multievidence.py}):

\begin{table*}[t]
\centering
\scriptsize
\resizebox{\textwidth}{!}{%
\begin{tabular}{lrrrrr}
\toprule
system & nDCG@10 & relevant docs in top-10 & questions with >=2 & step coverage of top-10 & steps retrievable @1000 \\
\midrule
dense & 0.2943 & 1.03 & 31.8\% & 51.3\% & 96.3\% \\
+ zero-shot CE & 0.3249 & 1.11 & 35.6\% & 54.9\% & 96.3\% \\
+ SFT CE (frozen) & 0.3581 & 1.21 & 38.3\% & 59.5\% & 96.3\% \\
\bottomrule
\end{tabular}}%
\end{table*}

Supervision improves every column but leaves most of the structure unused. At depth 1000 seven of ten steps are retrievable and still four of ten never reach the top ten; at the depth the reranker sees (100 candidates per query) only 80.7\% of annotated steps are reachable and 88.3\% of questions keep a positive, so a fifth of the step structure is lost before the cross-encoder runs. Truncating the fused list from 1000 to 100 rows leaves its top ten unchanged (59.45\%, \texttt{results/MULTIEVIDENCE\_sft\_e2bc\_sub100.json}), so depth moves reachability only; the 51.4\% is the cross-encoder’s own ordering over those 100 candidates. The retrieval floor is small but visible: in 10 of 519 dev questions (1.9\%) no relevant document appears anywhere in the dense top-1000, and 52 of the 138 questions the frozen system scores at zero (38\%) have none in the top 100, where no reranker can lift them. Coverage tracks nDCG@10 (Spearman 0.69; all-covered questions average 0.595 against 0.187, \texttt{results/STEP\_COVERAGE\_CORR.json}), though partly mechanically, so we read it as a check rather than a result.

The gain does not concentrate on the questions that need several evidence documents: on the 151 two-step dev questions it is +0.0200 (CI [-0.0033, +0.0463]) against +0.0711 on single-step ones (Appendix A.4).

\subsection*{5.3 Query truncation drops the conditions}
In 1,730 train+dev queries, 44.8\% (46.4\% of dev: 241 of 519) carry at least one condition (year 14.3\%, temporal relation 19.0\%, exclusion 16.3\%, comparison 14.0\%) that appears \textbf{only in the body}, so it never reaches the cross-encoder, which is trained on the title alone; the dense stage reads the full question. That truncation is not, by itself, what the ranker is missing: on the 241 condition-lost queries the frozen system scores \textbf{0.3656 against 0.3518} on the rest, and its gain over zero-shot CE is about twice as large there (+0.0481 vs +0.0246), although the dense stage already reads the full question and only the reranker does not. A same-budget body-visible composer therefore has to \emph{earn} the claim, so we pre-register that its gain must concentrate on the lost subset. Two 64-word composers were built - a cue-sentence selector (title plus only the body sentences carrying a condition) and the plain full question (title plus the first 64 body words). On the 45-query internal split no representation wins outside noise: title 0.2827, full 0.2869, selector 0.2744, fixed body 0.2703, frequency summary 0.2614 (\texttt{results/QVARIANT\_WINNER.json}); the \textasciitilde{}+-0.03 resolution of that split (S7) makes dev the decisive test, and the selector - the more faithful of the two - was the worst, so dev carries the \emph{generous} arm. Trained end to end on the 64-word full question it scores \textbf{0.3443} against 0.3581 (Delta -0.0138, query-clustered CI [-0.0272, -0.0005], domain-clustered [-0.0338, +0.0052]) and loses on both subsets: 0.3566 vs 0.3656 on the 241 condition-lost queries and 0.3296 vs 0.3518 on the other 278. The loss is thus \emph{smaller} where the arm was designed to help (-0.0090 vs -0.0222), consistent with body conditions not being the missing signal, but no subset gains, so the pre-registered bar is not met (\texttt{results/CONSTRAINT\_SUBSET\_cx\_qv.json}).

\subsection*{5.4 Temporal modelling did not help}
Ten probes - year/anchor features, a semantic time-consistency score, oracle and LLM step decomposition, HyDE, temporal hard negatives, a time-aware constraint rerank, a train-fitted gate, a fixed year-consistency reordering, and a coverage diagnostic. Exactly one was strictly pre-registered: it gained +0.0039 ([-0.0023, +0.0103], domain [+0.0006, +0.0142], P=0.893) against a +0.02 bar. The year rule fails cleanly: ordering the frozen candidates by year consistency lowers nDCG@10 in 7 of 8 anchored domains (macro 0.2980 to 0.2498). That is a fixed rule losing, not an upper bound - the reordered list is not the best possible one, and the rule itself is what a temporal ranker would have to beat. In this corpus a document mentioning another period is often still the right answer.

\subsection*{5.5 The rest of the ladder, in brief}
Three further steps are reported in full in Appendix A. \textbf{Loss forms}: averaging the log-loss over all positives (0.2809), an nDCG-weighted permutation loss (0.2764) and domain-sqrt sampling (0.2912) all sit within the 0.010 tie band of the current objective (0.2865) on the 45-query split and were not promoted to dev, so their dev effect is unmeasured. \textbf{Step supervision}: requiring every annotated step to have a positive ranked high gives 0.3593 and 0.3600 against 0.3581 (domain-clustered CIs [-0.0021, +0.0048] and [-0.0014, +0.0058]) - a wash, and coverage stays at 59.5\% / 59.7\% rather than rising. \textbf{Corrections}: restoring the 140 empty-pool queries is also a wash (0.3565, -0.0017 [-0.0088, +0.0061]); the 64-word full question is the one arm that pays (0.3443, -0.0138 [-0.0272, -0.0005]). With S5.3 and S5.4 the message is consistent: no retrained arm beats the frozen system, and only the query-rewriting arm loses to it.

\subsection*{5.6 A better cross-encoder is not a better system}
The checkpoint screen scores the cross-encoder alone; the system is the dense+CE fusion, and the two rank the arms differently. \texttt{cx\_qv} is the stronger cross-encoder by 0.0387 (0.3152 against 0.2765) and 0.0132 \emph{worse} fused at the frozen weight.

Three mechanisms explain it (519 dev queries, frozen dense top-1000, per-query z-normalised 1:1 fusion). \emph{(i) Anchoring}: the fused top ten is 73-75\% dense documents and only 54-60\% of it is also in the CE top ten, so a changed ordering must displace dense documents to matter. \emph{(ii) Attenuation}: a CE gain transfers at slope 0.29 - +0.2082 on the 169 queries the CE improves becomes +0.0438 fused, and -0.1661 on the 135 it hurts becomes -0.0665. \emph{(iii) Profile}: on the 215 queries where the CE’s own nDCG@10 does not move (|delta| <= 0.01) the fused score still falls 0.0113, and a transplant isolates the two - \texttt{cx\_qv}’s ordering with \texttt{lx\_possum}’s score profile scores 0.3500 (-0.0075) while the reverse falls to 0.3421 (-0.0154).

Re-weighting the fusion narrows the gap but does not close it. The weight is a free parameter and the optimum is model-specific: on the 45 held-out queries \texttt{lx\_possum} prefers 0.75 and \texttt{cx\_qv} 1.25, and carrying those to dev gives 0.3552 and 0.3509 - a 0.0023 loss for the model the weight was tuned for and a 0.0066 gain for the other, still 0.0072 under the frozen system. In the controlled probe, where the z-profile is computed over the same 100 candidates, the sweep settles at 1.0-1.5 for \texttt{lx\_possum} (0.3575, 0.3563) and at 1.5-2.0 for \texttt{cx\_qv} (0.3532, 0.3546), so the better standalone model stays the worse fused one at every weight we tried. A calibration-free RRF (k = 60) gives 0.3433 and 0.3462: it orders the two correctly and trails the frozen system by 0.012-0.015. On those same queries the model-level inversion reproduces with the candidate list fixed (\texttt{cx\_qv} 0.2893 alone against 0.2462, and 0.3197 against 0.3325 fused); within a model both criteria prefer epoch 1, and that preference is itself noise (epoch 1 on dev: 0.3560 against 0.3575, CI [-0.0047, +0.0020], 423 of 519 tied). A single-stage metric is thus a good ranker of cross-encoders and a poor ranker of the system they feed.

\section*{6 Reproducibility}
Every number in the tables is recomputed from released artefacts rather than transcribed: per-query score JSON for each arm, the dev TREC runs, and the official qrels. \texttt{scripts/check\_claims.py} recomputes the headline figures from those files and fails if one drifts, and the pipeline itself is three commands - dense retrieval at depth 1,000, reranking at depth 100, then a z-normalised weighted sum with the weights fixed at 1:1. The cost is one consumer GPU (an RTX 4060 Laptop with 8 GB): encoding the 1.17 M documents is the dominant term, about six hours at the 52 documents/s we measure; reranking 519 queries at depth 100 takes about 80 minutes; a two-epoch training run takes tens of minutes. The three seeds and the two checkpoint epochs per run are released as runs, so the retraining band in S4 can be re-derived without a GPU.

\section*{7 Limitations}
Domain-macro lets 3-query domains weigh as much as 240-query ones, so a single-domain regression can move the headline; the temporal findings are about \emph{textual} time only, not about systems that have timestamps; and the supervised results are single-corpus and single-family (one dense model, one cross-encoder). The window/pooling schedule (250-token windows, stride 192, max over windows, 256-token query cap) was fixed a priori and is not ablated; the dense encoder’s pooling was ablated late and is a wash (CLS 0.2956 against 0.2943 for mean); the 1:1 fusion weight and the two-epoch schedule were picked on the training-internal split, so the frozen number carries that selection. The step-supervision results are reported at one seed; a second seed and distillation of a larger reranker \citep{hinton2015distilling} were not affordable within the GPU budget.

Selection is noisy by construction: the 45-query internal split improves the sum-positive arm by +0.0135 with a domain-clustered interval of [-0.0317, +0.0571], i.e. it cannot resolve differences below about 0.04. Every arm here is therefore screened internally and judged on dev with a paired domain bootstrap (half-width 0.007-0.014); dev gains below +0.010 are treated as unresolved. The screen did not promote every loser (S5.5) - what it cannot do is rank these arms: the five it scored lie within 0.006 of each other, dev puts them between -0.0017 and +0.0019 of the frozen system, and every interval includes zero. Retraining is not free: fresh pools and seeds give 0.3575 / 0.3528 / 0.3527 and the current code path 0.3575, so 0.3581 sits at the top of a 0.3527-0.3581 band rather than being a property of the recipe.

\section*{8 Error analysis}
A quarter of dev questions (138/519, 26.6\%) score exactly zero. The split separates cleanly by candidate-pool recall: \textbf{52 (37.7\%) have no relevant document in the dense top-100} - no top-100 reranker can score them - and 86 (62.3\%) have the gold in the pool and lose it to ordering. Zeros cluster where the corpus is thin (history 46, politics 14, bitcoin 13). Three worked examples, one per mechanism (\texttt{results/DEV\_ERROR\_ANALYSIS.md}):

\begin{table}[t]
\centering
\footnotesize
\resizebox{\columnwidth}{!}{%
\begin{tabular}{lrrrr}
\toprule
failure & query & gold & dense rank & nDCG@10 \\
\midrule
pool miss & travel/201562\_1094 & 2 & 656, not retrieved & 0.0000 \\
ordering & history/70583\_271 & 1 & 24 (CE 14, fused 17) & 0.0000 \\
step coverage & history/59309\_621 & 2 & 1, 10 (CE 1, 72) & 0.6131 \\
\bottomrule
\end{tabular}}%
\end{table}

The pool miss is a PNR booking-limit question whose evidence sits outside the 100 candidates the reranker sees (dense rank 656); a deeper second stage could reach that document, but nothing downstream of the reranker can. The ordering case ("How do we know when Nero was born?") has its gold at dense rank 24, moves to 14 for the cross-encoder and 17 fused, and scores zero. The step case ("How many people watched the moon landings after Apollo 11?") keeps the Apollo 11 figures at rank 1 but pushes the comparison document from dense rank 10 to 72 for the title-only cross-encoder: one of two evidence documents survives, nDCG@10 stays at 0.61, and step coverage fails. That shape is common, not anomalous: only 38.3\% of questions place two or more relevant documents in the top ten (S5.2).

\section*{9 Conclusion}
A zero-shot dense stage plus a cross-encoder fine-tuned on the 1,071 questions with an in-pool positive gives +0.0332 domain-macro nDCG@10 over the same pipeline with the cross-encoder left at zero-shot (0.3249), on all 13 dev domains, stable across seeds (0.347-0.358). Two failures remain: an incomplete pool (1.9\% of questions have no gold in the dense top-1,000) and multi-evidence coverage (59.5\% of steps in the fused top ten). Every explicit temporal intervention - ten probes, one pre-registered test, step supervision, query reformulation, loss variants, backfilling, rank fusion - is neutral or negative. Most transferable: a cross-encoder 0.04 better on its own metric can be worse fused, because fusion anchors on the first stage and re-normalises the second’s scores.

\begin{thebibliography}{99}
\small
\bibitem[Abdallah et al.(2026)]{abdallah2026tempo} Abdelrahman Abdallah, Mohammed Ali, Muhammad Abdul-Mageed, and Adam Jatowt. 2026. TEMPO: A Realistic Multi-Domain Benchmark for Temporal Reasoning-Intensive Retrieval. arXiv:2601.09523.
\bibitem[Abdallah et al.(2026)]{abdallah2026reteco} Abdelrahman Abdallah, Mohammed Ali, Muhammad Abdul-Mageed, Kevin Duh, and Adam Jatowt. 2026. RETECO: A SemEval-2027 Shared Task on Reasoning-Oriented Retrieval. Task proposal, University of Innsbruck. The organizers ask participants to cite TEMPO (and RECOR for Track 2) until the task description paper appears.
\bibitem[Robertson(2009)]{robertson2009bm25} Stephen Robertson and Hugo Zaragoza. 2009. The Probabilistic Relevance Framework: BM25 and Beyond. \emph{Foundations and Trends in Information Retrieval}, 3(4):333-389.
\bibitem[Nogueira(2019)]{nogueira2019passage} Rodrigo Nogueira and Kyunghyun Cho. 2019. Passage Re-ranking with BERT. arXiv:1901.04085.
\bibitem[Cormack et al.(2009)]{cormack2009rrf} Gordon V. Cormack, Charles L. A. Clarke, and Stefan Buettcher. 2009. Reciprocal Rank Fusion Outperforms Condorcet and Individual Rank Learning Methods. In \emph{SIGIR 2009}, pages 758-759.
\bibitem[MSR-TR-2010-82(2010)]{burges2010lambdarank} Christopher J. C. Burges. 2010. From RankNet to LambdaRank to LambdaMART: An Overview. Technical Report MSR-TR-2010-82, Microsoft Research.
\bibitem[Wang et al.(2018)]{wang2018lambdaloss} Xuanhui Wang, Cheng Li, Nadav Golbandi, Michael Bendersky, and Marc Najork. 2018. The LambdaLoss Framework for Ranking Metric Optimization. In \emph{CIKM 2018}, pages 1313-1322.
\bibitem[Jarvelin(2002)]{jarvelin2002ndcg} Kalervo Jarvelin and Jaana Kekalainen. 2002. Cumulated Gain-based Evaluation of IR Techniques. \emph{ACM Transactions on Information Systems}, 20(4):422-446.
\bibitem[Hinton et al.(2015)]{hinton2015distilling} Geoffrey Hinton, Oriol Vinyals, and Jeff Dean. 2015. Distilling the Knowledge in a Neural Network. arXiv:1503.02531.
\bibitem[Chen et al.(2024)]{chen2024m3} Jianlv Chen, Shitao Xiao, Peitian Zhang, Kun Luo, Defu Lian, and Zheng Liu. 2024. BGE M3-Embedding: Multi-Lingual, Multi-Functionality, Multi-Granularity Text Embeddings Through Self-Knowledge Distillation. arXiv:2402.03216.
\bibitem[card(2024)]{mxbai2024card} mixedbread.ai. 2024. mxbai-embed-large-v1 model card. https://huggingface.co/mixedbread-ai/mxbai-embed-large-v1 No peer-reviewed description exists for this checkpoint; we cite the model card and report the exact revision in \texttt{results/MODEL\_MANIFEST.json}.
\bibitem[Zhang et al.(2022)]{zhang2022hlatr} Yanzhao Zhang, Dingkun Long, Guangwei Xu, and Pengjun Xie. 2022. HLATR: Enhance Multi-Stage Text Retrieval with Hybrid List Aware Transformer Reranking. arXiv:2205.10569.
\bibitem[Min et al.(2021)]{mao2021jpr} Sewon Min, Kenton Lee, Ming-Wei Chang, Kristina Toutanova, and Hannaneh Hajishirzi. 2021. Joint Passage Ranking for Diverse Multi-Answer Retrieval. In \emph{EMNLP 2021}, pages 6997-7008.
\end{thebibliography}

\section*{Appendix A. The six-step ladder in full}
\subsection*{A.1 The loss form is not the bottleneck}
Three alternatives suggested by the coverage gap - averaging the log-loss over all positives, an nDCG-weighted permutation loss (LambdaLoss-style, cutoff 10), and domain-balanced (sqrt) sampling - were trained on the same 942 groups with the same hyper-parameters and scored on the same 45-query split:

\begin{table}[t]
\centering
\footnotesize
\resizebox{\columnwidth}{!}{%
\begin{tabular}{lrr}
\toprule
objective / sampler & best epoch & val domain-macro nDCG@10 \\
\midrule
current (log-sum-exp over positives) & 2 & 0.2865 \\
mean over positives & 2 & 0.2809 \\
nDCG-weighted permutations & 2 & 0.2764 \\
current + sqrt(domain) sampling & 2 & 0.2912 \\
zero-shot reference & - & 0.2730 \\
\bottomrule
\end{tabular}}%
\end{table}

Selecting on domain-macro rather than the plain mean picks the same epoch in all four arms, so the official aggregation does not move the checkpoint; the split cannot resolve the gaps (+-0.03, S7), so only sqrt sampling’s +0.0047 ([-0.0164, +0.0285]) is even measurable, and the pre-registered 0.010 tie band keeps the current objective. Sqrt sampling’s one visible effect is raising the share of queries with a relevant document in the top ten from 93.3\% to 95.6\% without moving nDCG: it makes the easy decision easier.

Dev measures the incumbent itself. Retrained through the current code path (same 942 groups, two epochs, its own val-preferred epoch) it scores \textbf{0.3575}, 0.0006 below the frozen artifact (domain-clustered CI [-0.0037, +0.0021]), so re-running the recipe is free within measurement error and the frozen artifact is a draw from the same band. The three replacement loss forms stay unmeasured on dev - the screen separates them by at most 0.015, below the dev half-width of 0.007-0.014 - so dev cannot say they lose; what it can say is that the screen ranks arms at 0.2764-0.2912, cannot resolve 0.003, and that the arm it did put on dev reproduces the frozen number.

\subsection*{A.2 Forcing every evidence step into the top ten does not pay}
If the fused top ten covers only 59.5\% of annotated steps (51.4\% for the cross-encoder alone), the mechanical remedy is to require every step to have a positive ranked high. We trained the step-augmented objective of S3 at two weights and scored the systems on dev with the frozen fusion:

\begin{table*}[t]
\centering
\scriptsize
\resizebox{\textwidth}{!}{%
\begin{tabular}{lrrrrr}
\toprule
system & dev domain-macro & Delta vs frozen & domain-clustered 95\% CI & steps in top 10 & questions with 2+ relevant in top 10 \\
\midrule
frozen (no step term) & 0.3581 & - & - & 59.45\% & 38.31\% \\
L\_step, lambda = 0.3 & 0.3593 & +0.0012 & [-0.0021, +0.0048] & 59.54\% & 38.20\% \\
L\_step, lambda = 1.0 & 0.3600 & +0.0019 & [-0.0014, +0.0058] & 59.69\% & 38.47\% \\
\bottomrule
\end{tabular}}%
\end{table*}

The coverage column counts the fused list, whose top ten does not change when the list is truncated from 1000 to 100 rows (59.45\%); the cross-encoder’s own ordering over the same candidates covers 51.39\% before the term and 48.14\% / 45.78\% with it, so the degradation is a property of the ordering, not of depth.

The term does not pay and does not move the quantity it optimises: coverage stays at 59.5\% / 59.7\%. One reading fits: a step’s positives are mostly documents the ordinary listwise term already covers (3.44 relevant documents over 2.34 steps, so a step rarely adds a \emph{different} document), and the log-sum-exp inside a step is dominated by the step whose positive is hardest to retrieve, so the gradient chases documents that are absent or buried. The annotations are not at fault: every one of the 2,873 step-level positives is also a query-level positive. The instructive part is the selection: on the internal split lambda = 1.0 was the \emph{best} of six arms trained that day (0.2913, above 0.2865), and on dev the two step arms are indistinguishable from the frozen system - which is why S3 fixes the fusion form in advance and S7 refuses to select on 45 questions.

\subsection*{A.3 Two further corrections that also lose}
The frozen system feeds the reranker a title-only query and excludes the 140 training queries whose dense top-100 contains no positive. Both were corrected with the frozen recipe and scored on dev with the frozen fusion:

\begin{table}[t]
\centering
\footnotesize
\resizebox{\columnwidth}{!}{%
\begin{tabular}{lrrrr}
\toprule
system & dev domain-macro & Delta vs frozen & query-clustered 95\% CI & domain-clustered 95\% CI \\
\midrule
frozen & 0.3581 & - & - & - \\
full-question query, 64-word cap & 0.3443 & -0.0138 & [-0.0272, -0.0005] & [-0.0338, +0.0052] \\
+ 140 backfilled queries, 294 positives & 0.3565 & -0.0017 & [-0.0088, +0.0061] & [-0.0061, +0.0061] \\
\bottomrule
\end{tabular}}%
\end{table}

The two no longer cost the same: the backfill is a wash and the query rewrite pays. For the query arm the loss is smaller where the body conditions live (S5.3), so the extra text carries signal but not enough to pay for the harm on the other 278 queries. For the backfill arm the 294 restored positives are, by construction, documents the first stage never retrieves, so training on them spends capacity on a decision never made at inference, where the candidate list is the dense top-100 - and on dev that costs nothing measurable. The unchanged objective re-run through the current code path scores 0.3575, 0.0006 below the frozen artifact, so the full-question arm loses 0.0132 against its own baseline and the step arms sit 0.0018 / 0.0025 above it. Dev therefore supports a deployment decision - six of the ladder’s steps retrained, none beating 0.3581, and the sixth (dense pooling) a wash at +0.0014 - plus a negative result against a same-code baseline that the 45-query screen cannot rank at all.

\subsection*{A.4 Where the gain sits: not on the multi-step questions}
Partitioning dev by the number of official evidence steps a question needs (\texttt{results/STEP\_SPLIT.md}), the frozen system’s gain over zero-shot CE is largest on single-step questions and weakest on two-step ones:

\begin{table*}[t]
\centering
\scriptsize
\resizebox{\textwidth}{!}{%
\begin{tabular}{lrrrrr}
\toprule
partition & queries & zero-shot CE & frozen SFT & delta & 95\% CI (domain-clustered) \\
\midrule
1 step & 75 & 0.3185 & 0.3896 & +0.0711 & [+0.0244, +0.1056] \\
2 steps & 151 & 0.3857 & 0.4058 & +0.0200 & [-0.0033, +0.0463] \\
3+ steps & 258 & 0.3196 & 0.3511 & +0.0315 & [+0.0084, +0.0571] \\
none & 35 & 0.1579 & 0.1705 & +0.0126 & [+0.0030, +0.0524] \\
\bottomrule
\end{tabular}}%
\end{table*}

The objective buys a better \emph{single-document} decision almost everywhere but does not teach evidence assembly: on the 151 two-step questions the gain is within noise. That is the gap the step term targets, and it fixes the bar those arms have to clear.

\end{document}
